%% abtex2-modelo-trabalho-academico.tex, v-1.9.7 laurocesar
%% Copyright 2012-2018 by abnTeX2 group at http://www.abntex.net.br/ 
%%
%% This work may be distributed and/or modified under the
%% conditions of the LaTeX Project Public License, either version 1.3
%% of this license or (at your option) any later version.
%% The latest version of this license is in
%%   http://www.latex-project.org/lppl.txt
%% and version 1.3 or later is part of all distributions of LaTeX
%% version 2005/12/01 or later.
%%
%% This work has the LPPL maintenance status `maintained'.
%% 
%% The Current Maintainer of this work is the abnTeX2 team, led
%% by Lauro César Araujo. Further information are available on 
%% http://www.abntex.net.br/
%%
%% This work consists of the files abntex2-modelo-trabalho-academico.tex,
%% abntex2-modelo-include-comandos and abntex2-modelo-references.bib
%%

% ------------------------------------------------------------------------
% ------------------------------------------------------------------------
% abnTeX2: Modelo de Trabalho Academico (tese de doutorado, dissertacao de
% mestrado e trabalhos monograficos em geral) em conformidade com 
% ABNT NBR 14724:2011: Informacao e documentacao - Trabalhos academicos -
% Apresentacao
% ------------------------------------------------------------------------
% ------------------------------------------------------------------------

\documentclass[
	% -- opções da classe memoir --
	12pt,						% tamanho da fonte
	openright,					% capítulos começam em pág ímpar (insere página vazia caso preciso)
	oneside,					% para impressão em recto e verso. Oposto a oneside
	a4paper,					% tamanho do papel. 
	% -- opções da classe abntex2 --
	%chapter=TITLE,				% títulos de capítulos convertidos em letras maiúsculas
	%section=TITLE,				% títulos de seções convertidos em letras maiúsculas
	%subsection=TITLE,			% títulos de subseções convertidos em letras maiúsculas
	%subsubsection=TITLE,		% títulos de subsubseções convertidos em letras maiúsculas
	% -- opções do pacote babel --
	english,					% idioma adicional para hifenização
	french,						% idioma adicional para hifenização
	spanish,					% idioma adicional para hifenização
	brazil						% o último idioma é o principal do documento
	]{abntex2}


\usepackage{lmodern}			% Usa a fonte Latin Modern			
\usepackage[T1]{fontenc}		% Selecao de codigos de fonte.
\usepackage[utf8]{inputenc}		% Codificacao do documento (conversão automática dos acentos)
\usepackage{indentfirst}		% Indenta o primeiro parágrafo de cada seção.
\usepackage{color}				% Controle das cores
\usepackage{graphicx}			% Inclusão de gráficos
\usepackage{microtype} 			% para melhorias de justificação
\usepackage{lipsum}				% para geração de dummy text
\usepackage[brazilian,hyperpageref]{backref}	 % Paginas com as citações na bibl
\usepackage[alf]{abntex2cite}	% Citações padrão ABNT
\usepackage{multirow}

\usepackage{pdfpages}
\usepackage{forest}
\usepackage{fontawesome}

\definecolor{thered}{rgb}{0.65,0.04,0.07}
\definecolor{thegreen}{rgb}{0.06,0.44,0.08}
\definecolor{thegrey}{gray}{0.5}
\definecolor{theshade}{rgb}{1,1,0.97}
\definecolor{theframe}{gray}{0.6}
\definecolor{blue}{RGB}{41,5,195} 

\usepackage{listings}
\lstset{
	basicstyle=\small, % print whole listing small
	keywordstyle=\color{blue}\bfseries, %\underbar,% underlined bold black keywords
	% keywords=[0]{uint32_t, exit},
	% keywordstyle=[0]\color{thered},
	identifierstyle=, % nothing happens
	commentstyle=\color{gray!85}, % white comments
	stringstyle=\color{thered}\ttfamily, % typewriter type for strings
	showstringspaces=false, % no special string spaces
	numbers=left,
	numberstyle=\tiny,
	stepnumber=2
}

\definecolor{folderbg}{RGB}{124,166,198}
\definecolor{folderborder}{RGB}{110,144,169}

\def\Size{4pt}
\tikzset{
  folder/.pic={
    \filldraw[draw=folderborder,top color=folderbg!50,bottom color=folderbg]
      (-1.05*\Size,0.2\Size+5pt) rectangle ++ (.75*\Size,-0.2\Size-5pt);  
    \filldraw[draw=folderborder,top color=folderbg!50,bottom color=folderbg]
      (-1.15*\Size,-\Size) rectangle (1.15*\Size,\Size);
  }
}


% Configurações do pacote backref
% Usado sem a opção hyperpageref de backref
\renewcommand{\backrefpagesname}{Citado na (s) página (s):~}
% Texto padrão antes do número das páginas
\renewcommand{\backref}{}
% Define os textos da citação
\renewcommand*{\backrefalt}[4]{
	\ifcase#1
		Nenhuma citação no texto.%
	\or%
		Citado na página #2.%
	\else
		Citado #1 vezes nas páginas #2.%
	\fi}%




%-------------------------------------------------------------------------------------
%		DADOS PESSOAIS
%		Informações básicas sobre os trabalho e os autores 
%-------------------------------------------------------------------------------------
\titulo{Título da Dissertação}
\autor{Nome do Autor}
\local{Salvador}
\data{Dezembro 2022}
\orientador{Prof. Dr. Nome do Orientador}
\coorientador{Prof. Dr. Nome do Coorientador}
\tipotrabalho{Dissertação (Mestrado)} 
\preambulo{Dissertação de Mestrado apresentada à Comissão Acadêmica Institucional do PROFMAT-UFBA como requisito parcial para obtenção do título de Mestre em Matemática.}

\instituicao{%
  Universidade Federal da Bahia - UFBA
  \par
  Instituto de Matemática - IM
  \par
  Sociedade Brasileira de Matemática - SBM
}

%	Configurações de aparência do PDF final
\input{config_pdf}





%%%%%%%%%%%%%%%%%%%%%%%%%%%%%%%%%%%%%%%%%%%%%%%%%%%%%%%%%%%%%%%%%%%%%%%%%%%%%%%%%%%%%%
%																					 %
% 							 Início do documento								     %
%																					 %
%%%%%%%%%%%%%%%%%%%%%%%%%%%%%%%%%%%%%%%%%%%%%%%%%%%%%%%%%%%%%%%%%%%%%%%%%%%%%%%%%%%%%%
\begin{document}

%\selectlanguage{english}
\selectlanguage{brazil}

% Retira espaço extra obsoleto entre as frases.
\frenchspacing 

%-------------------------------------------------------------------------------------
% 		ELEMENTOS PRÉ-TEXTUAIS
%		Insire elementos pré textuais, ficha catalográfica 
%-------------------------------------------------------------------------------------
% \pretextual

\imprimircapa

\imprimirfolhaderosto*	% (o * indica que haverá a ficha bibliográfica)


\include{pre-textuais/ficha_catalografica}
\include{pre-textuais/errata}
\include{pre-textuais/folha_aprovacao}
\include{pre-textuais/dedicatoria}
\include{pre-textuais/agradecimentos}
\include{pre-textuais/epigrafe}
\include{pre-textuais/resumo}
\include{pre-textuais/abstract}
%\include{pre-textuais/listas_1}
%\include{pre-textuais/listas_siglas}
%\include{pre-textuais/listas_simbolos}

% inserir o sumario
\pdfbookmark[0]{\contentsname}{toc}
\tableofcontents*
\cleardoublepage%

%-------------------------------------------------------------------------------------
% 		ELEMENTOS TEXTUAIS
%		Informações relevantes ao trabalho, capítulos e suas respectivas seções.
%-------------------------------------------------------------------------------------
\textual%

\include{capitulos/capitulo_introducao}

% PARTE
% ----------------------------------------------------------
\part{Preparação da pesquisa}
\include{capitulos/capitulo_01}

% PARTE
% ----------------------------------------------------------                                  
\part{Referenciais teóricos}
\include{capitulos/capitulo_02}

% PARTE
% ----------------------------------------------------------  
\part{Resultados}
\include{capitulos/capitulo_03}
\include{capitulos/capitulo_04}

% ----------------------------------------------------------
% Finaliza a parte no bookmark do PDF
% para que se inicie o bookmark na raiz
% e adiciona espaço de parte no Sumário
% ----------------------------------------------------------
\phantompart%

% Conclusão
% ---
\include{capitulos/capitulo_05}

%-------------------------------------------------------------------------------------
%		CONTEÚDO PÓS-TEXTUAL
%		Parte do documento que aparece após o capítulo de conclusão, incluindo 
%		bibliografia e anexos.
%-------------------------------------------------------------------------------------
\postextual\bibliography{references} 

%\glossary

% Apendices
%-------------------------------
% \begin{apendicesenv}
% \partapendices
% \include{apendices/apendiceA}
% \end{apendicesenv}

% Anexos
%-------------------------------
% \begin{anexosenv}
% \partanexos
% \include{anexos/anexoA}
% \include{anexos/anexoB}
% \include{anexos/anexoC}
% \end{anexosenv}

% Indice remissivo
%-------------------------------
% \phantompart
% \printindex

\end{document}
